\section{2026-05-14}

\subsection{Smooth Structures and Smooth Manifolds}

\begin{definition}
Let $M$ be a topological manifold. A \vocab{smooth structure} on $M$ is a maximal smooth atlas $\mathcal A$ on $M$.
\end{definition}

\begin{definition}
A \vocab{smooth manifold} is a pair
\[
  (M^n,\mathcal A),
\]
where $M^n$ is a topological $n$-manifold and $\mathcal A$ is a smooth structure on $M$.
\end{definition}

Usually we use the shorthand notation
\[
  M=(M^n,\mathcal A).
\]
If
\[
  (U,\varphi)\in \mathcal A,
\]
then $(U,\varphi)$ is called a smooth chart of $M$.

\begin{remark}
The same construction can be modified by replacing ``smooth'' with other regularity classes. For example, one can define
\[
  C^k\text{-manifolds},\qquad C^\omega\text{-manifolds},
\]
where $C^\omega$ means real analytic.
\end{remark}

\subsection{Local Properties of Smooth Charts}

\begin{exercise}
Let $(U,\varphi)\in\mathcal A$, and let $U'\subset U$ be open. Then
\[
  (U',\varphi|_{U'})\in \mathcal A.
\]
\end{exercise}

\begin{proof}
Since $U'\subset U$ is open, $\varphi(U')$ is open in $\RR^n$, and
\[
  \varphi|_{U'}:U'\to \varphi(U')
\]
is a homeomorphism. Thus $(U',\varphi|_{U'})$ is a chart.

It is smoothly compatible with every chart in $\mathcal A$, because its transition maps are restrictions of the corresponding transition maps for $(U,\varphi)$. Since $\mathcal A$ is maximal, it follows that
\[
  (U',\varphi|_{U'})\in \mathcal A.
\]
\end{proof}

\begin{exercise}
Let $(U,\varphi)\in\mathcal A$, and let
\[
  \chi:\varphi(U)\to \chi(\varphi(U))\subset \RR^n
\]
be a diffeomorphism onto an open subset of $\RR^n$. Then
\[
  (U,\chi\circ\varphi)\in \mathcal A.
\]
\end{exercise}

\begin{proof}
The map
\[
  \chi\circ\varphi:U\to \chi(\varphi(U))
\]
is a homeomorphism, hence $(U,\chi\circ\varphi)$ is a chart.

For any $(V,\psi)\in\mathcal A$, the transition map
\[
  \psi\circ(\chi\circ\varphi)^{-1}
  =
  \psi\circ\varphi^{-1}\circ\chi^{-1}
\]
is smooth, because it is a composition of smooth maps. Similarly,
\[
  (\chi\circ\varphi)\circ\psi^{-1}
  =
  \chi\circ\varphi\circ\psi^{-1}
\]
is smooth. Hence $(U,\chi\circ\varphi)$ is smoothly compatible with all charts in $\mathcal A$. By maximality,
\[
  (U,\chi\circ\varphi)\in\mathcal A.
\]
\end{proof}

\begin{exercise}
Let $\varphi:U\to \RR^n$ be an injective map, where $U\subset M$ is open. Suppose that for every $p\in U$, there exists an open neighborhood $U_p\subset U$ of $p$ such that
\[
  (U_p,\varphi|_{U_p})\in\mathcal A.
\]
Then
\[
  (U,\varphi)\in\mathcal A.
\]
\end{exercise}

\begin{proof}
The condition is local. Since each $\varphi|_{U_p}$ is a coordinate map, the image $\varphi(U_p)$ is open in $\RR^n$. Because
\[
  \varphi(U)=\bigcup_{p\in U}\varphi(U_p),
\]
the image $\varphi(U)$ is open in $\RR^n$.

Moreover, $\varphi$ is a homeomorphism onto its image because this can be checked locally on the open cover $\{U_p\}_{p\in U}$. Thus $(U,\varphi)$ is a chart.

It remains to check smooth compatibility. Let $(V,\psi)\in\mathcal A$. For each $p\in U\cap V$, choose $U_p$ as above. On
\[
  \varphi(U_p\cap V),
\]
the transition map
\[
  \psi\circ\varphi^{-1}
\]
is equal to
\[
  \psi\circ(\varphi|_{U_p})^{-1},
\]
which is smooth because $(U_p,\varphi|_{U_p})\in\mathcal A$. Hence $\psi\circ\varphi^{-1}$ is smooth locally, and therefore smooth. The other transition map is treated similarly.

Thus $(U,\varphi)$ is smoothly compatible with all charts in $\mathcal A$. By maximality,
\[
  (U,\varphi)\in\mathcal A.
\]
\end{proof}

\begin{example}
Let $M=\RR$, and let $\mathcal A$ be the standard smooth structure, namely the maximal smooth atlas containing
\[
  (\RR,\id_{\RR}).
\]
Let $\mathcal A^*$ be the maximal smooth atlas containing
\[
  (\RR,\varphi),
  \qquad
  \varphi(x)=x^3.
\]
Then
\[
  \mathcal A^*\neq \mathcal A.
\]
Indeed, the two charts $(\RR,\id_{\RR})$ and $(\RR,\varphi)$ are not smoothly compatible, because one transition map is
\[
  \id_{\RR}\circ\varphi^{-1}(x)=\sqrt[3]{x},
\]
which is not smooth at $0$.

However, the smooth manifolds
\[
  (\RR,\mathcal A)
  \qquad\text{and}\qquad
  (\RR,\mathcal A^*)
\]
are diffeomorphic. For example, the map
\[
  \varphi:(\RR,\mathcal A^*)\to(\RR,\mathcal A),
  \qquad
  x\mapsto x^3,
\]
is a diffeomorphism.
\end{example}

\begin{remark}
This example shows that two smooth structures on the same underlying topological manifold may be different, even though the resulting smooth manifolds can still be diffeomorphic.
\end{remark}

\subsection{Standard Examples and Open Submanifolds}

\begin{example}
The standard smooth structure on $\RR^n$ is the maximal smooth atlas containing
\[
  (\RR^n,\id_{\RR^n}).
\]
\end{example}

\begin{theorem}
Let $(M,\mathcal A)$ be a smooth manifold, and let
\[
  M'\subset M
\]
be open. Define
\[
  \mathcal A'
  :=
  \{(U,\varphi)\in\mathcal A\mid U\subset M'\}.
\]
Then
\[
  (M',\mathcal A')
\]
is also a smooth manifold.
\end{theorem}

\begin{proof}
Since $M'\subset M$ is open, it is a topological manifold with the subspace topology.

The charts in $\mathcal A'$ cover $M'$. Indeed, if $p\in M'$, choose a chart $(U,\varphi)\in\mathcal A$ with $p\in U$. Then
\[
  U\cap M'
\]
is open in $U$, so by the restriction property,
\[
  (U\cap M',\varphi|_{U\cap M'})\in\mathcal A.
\]
Moreover its domain is contained in $M'$, hence it belongs to $\mathcal A'$.

Smooth compatibility is inherited from $\mathcal A$. Therefore $\mathcal A'$ is a smooth atlas on $M'$. Taking the maximal smooth atlas determined by $\mathcal A'$, we obtain the induced smooth structure on $M'$.
\end{proof}

\begin{example}
Let $S^n$ be the unit sphere. Its standard smooth structure is the maximal smooth atlas containing the stereographic projection charts
\[
  (U_N,\varphi_N),
  \qquad
  (U_S,\varphi_S),
\]
where $U_N=S^n\setminus\{N\}$ and $U_S=S^n\setminus\{S\}$.

Equivalently, it is the maximal smooth atlas containing the usual graph charts
\[
  (U_i^\pm,\varphi_i^\pm).
\]
\end{example}

\begin{example}
Let $V$ be an $n$-dimensional real vector space. Define
\[
  \mathcal A'
  :=
  \left\{
    (V,\varphi)
    \mid
    \varphi:V\to \RR^n \text{ is a vector space isomorphism}
  \right\}.
\]
Then $\mathcal A'$ is a smooth atlas, because the transition maps are linear isomorphisms
\[
  \RR^n\to \RR^n,
\]
hence smooth. The standard smooth structure on $V$ is the maximal smooth atlas containing $\mathcal A'$.
\end{example}

\begin{remark}
In general, if $\mathcal A$ is the maximal smooth atlas containing a smooth atlas $\mathcal A'$, and if a chart $(U,\varphi)$ is smoothly compatible with every chart in $\mathcal A'$, then
\[
  \mathcal A'\cup\{(U,\varphi)\}
\]
is still a smooth atlas. Hence
\[
  (U,\varphi)\in\mathcal A.
\]
\end{remark}

\begin{remark}
Reading examples include
\[
  S^n,\qquad \mathbb{RP}^n,\qquad \operatorname{GL}(n,\RR),
\]
and product manifolds.
\end{remark}

\begin{remark}
Digression.

Does every topological manifold have a smooth structure?

\[
  n\le 3:\quad \text{Yes, and unique up to diffeomorphism;}
\]
\[
  n\ge 4:\quad \text{No, and smooth structure may not be unique;}
\]

\end{remark}

\begin{remark}
Subquestion.

Are there exotic smooth structures on $S^n$?

Here ``exotic'' means not diffeomorphic to the standard one.

\[
  n\le 3:\quad \text{no;}
\]
\[
  n=4:\quad \text{unknown; \textit{Smooth Poincaré Conjecture}}
\]
\[
  n\ge 5:\quad \text{almost completely understood.} 
\]
On spheres are related to the group of homotopy spheres $\Theta_n$ which can be characterized using algebraic topology. case $n=126$ is often noted as unclear in this context.
\end{remark}

\subsection{Smooth Manifold Construction Lemma}

\begin{theorem}[Smooth Manifold Construction Lemma]
Let $M$ be a set, and let
\[
  \{\varphi_\alpha:U_\alpha\to \RR^n\}_{\alpha\in I}
\]
be a collection of injective maps, where $U_\alpha\subset M$. Suppose the following conditions hold.

\begin{enumerate}
  \item For every $\alpha,\beta\in I$,
  \[
    \varphi_\alpha(U_\alpha\cap U_\beta)
  \]
  is open in $\RR^n$.

  \item For every $\alpha,\beta\in I$, the transition map
  \[
    \varphi_\alpha\circ\varphi_\beta^{-1}
    :
    \varphi_\beta(U_\alpha\cap U_\beta)
    \to
    \varphi_\alpha(U_\alpha\cap U_\beta)
  \]
  is smooth.

  \item The set $M$ is covered by countably many of the sets $U_\alpha$.

  \item For any $p,q\in M$ with $p\neq q$, either there exists some $\alpha\in I$ such that
  \[
    p,q\in U_\alpha,
  \]
  or there exist $\alpha,\beta\in I$ such that
  \[
    p\in U_\alpha,\qquad q\in U_\beta,\qquad U_\alpha\cap U_\beta=\varnothing.
  \]
\end{enumerate}

Then $M$ has a unique topology and a unique smooth structure such that all pairs
\[
  (U_\alpha,\varphi_\alpha)
\]
are smooth charts.
\end{theorem}

\begin{proof}
Define a topology on $M$ by declaring that a subset $A\subset M$ is open if and only if
\[
  \varphi_\alpha(A\cap U_\alpha)
\]
is open in $\RR^n$ for every $\alpha\in I$.

With this topology, each $U_\alpha$ is open in $M$, and each
\[
  \varphi_\alpha:U_\alpha\to \varphi_\alpha(U_\alpha)
\]
is a homeomorphism onto an open subset of $\RR^n$.

The first two assumptions imply that the charts $(U_\alpha,\varphi_\alpha)$ are smoothly compatible. The countability assumption gives second-countability. The separation assumption gives the Hausdorff property. Therefore $M$ becomes a topological $n$-manifold.

Finally, let $\mathcal A$ be the maximal smooth atlas containing all charts
\[
  (U_\alpha,\varphi_\alpha).
\]
This gives the desired smooth structure. Uniqueness follows because the topology and the maximal smooth atlas are forced by the given charts.
\end{proof}

\subsection{Grassmannians}

\begin{example}[Grassmannians]
Let
\[
  1\le k\le n.
\]
The Grassmannian is
\[
  \operatorname{Gr}_k(\RR^n)
  :=
  \{V\subset \RR^n\mid V \text{ is a linear subspace and } \dim V=k\}.
\]
In particular,
\[
  \operatorname{Gr}_1(\RR^n)=\mathbb{RP}^{n-1}.
\]
\end{example}

Let
\[
  I
  :=
  \left\{
    (P,Q)
    \ \middle|\
    \begin{array}{l}
      P,Q\subset \RR^n \text{ are linear subspaces,}\\
      \RR^n=P\oplus Q,\ \dim P=k,\ \dim Q=n-k
    \end{array}
  \right\}.
\]
For
\[
  \alpha=(P,Q)\in I,
\]
define
\[
  U_\alpha
  :=
  \{V\in \operatorname{Gr}_k(\RR^n)\mid V\cap Q=\{0\}\}.
\]

\begin{lemma}
For every $V\in U_\alpha$, there exists a unique linear map
\[
  A_{P,Q,V}:P\to Q
\]
such that
\[
  V=\{x+A_{P,Q,V}x\mid x\in P\}.
\]
\end{lemma}

\begin{proof}
Since
\[
  \RR^n=P\oplus Q,
\]
every vector in $\RR^n$ can be written uniquely as $x+y$ with $x\in P$ and $y\in Q$.

If $V\in U_\alpha$, then
\[
  V\cap Q=\{0\}.
\]
The projection
\[
  \pi_P:V\to P
\]
is injective, because if $\pi_P(v)=0$, then $v\in Q$, hence $v\in V\cap Q=\{0\}$. Since
\[
  \dim V=\dim P=k,
\]
the map $\pi_P:V\to P$ is an isomorphism.

Therefore, for each $x\in P$, there is a unique vector $v\in V$ with $\pi_P(v)=x$. Write
\[
  v=x+A_{P,Q,V}x
\]
with $A_{P,Q,V}x\in Q$. This defines a unique linear map
\[
  A_{P,Q,V}:P\to Q.
\]
Hence
\[
  V=\{x+A_{P,Q,V}x\mid x\in P\}.
\]
\end{proof}

Define
\[
  \varphi_\alpha:U_\alpha\to L(P,Q)
\]
by
\[
  \varphi_\alpha(V)=A_{P,Q,V}.
\]
Since
\[
  L(P,Q)\cong \RR^{k(n-k)}
\]
after choosing bases of $P$ and $Q$, this gives a coordinate chart on the Grassmannian.

\subsection{Manifolds with Boundary}

\begin{definition}
Let
\[
  \mathbb H^n
  :=
  \{x=(x^1,\dots,x^n)\in\RR^n\mid x^n\ge 0\}
\]
be the closed upper half-space. Its boundary is
\[
  \partial \mathbb H^n
  :=
  \{x\in\mathbb H^n\mid x^n=0\}.
\]
\end{definition}

\begin{definition}
Let $M$ be an $n$-manifold with boundary. A point $p\in M$ is called a \vocab{boundary point} if there exists a chart
\[
  (U,\varphi)
\]
such that
\[
  p\in U
  \qquad\text{and}\qquad
  \varphi(p)\in \partial\mathbb H^n.
\]
The set of all boundary points is denoted by
\[
  \partial M.
\]
\end{definition}

For manifolds with boundary, transition maps are maps between open subsets of $\mathbb H^n$. Their smoothness means that all partial derivatives exist and extend smoothly up to the boundary. Equivalently, locally they extend to smooth maps on open subsets of $\RR^n$.

\begin{lemma}[Invariance of boundary]
Let $A,B\subset \mathbb H^n$ be open subsets in the subspace topology, and let
\[
  F:A\to B
\]
be a diffeomorphism. Then
\[
  F(A\cap \partial \mathbb H^n)=B\cap \partial \mathbb H^n.
\]
\end{lemma}

\begin{proof}
It suffices to prove that
\[
  F(A\cap \partial\mathbb H^n)\subset B\cap \partial\mathbb H^n,
\]
because the reverse inclusion follows by applying the same argument to
$F^{-1}$.

Suppose, for contradiction, that there exists
\[
  x\in A\cap \partial\mathbb H^n
\]
such that
\[
  y:=F(x)\in \operatorname{Int}\mathbb H^n.
\]
Let
\[
  G:=F^{-1}.
\]
Since $y$ is an interior point of $\mathbb H^n$, there is an ordinary open
neighborhood $W\subset \mathbb R^n$ of $y$ such that
\[
  W\subset B\cap \operatorname{Int}\mathbb H^n.
\]
Write
\[
  G=(G^1,\ldots,G^n).
\]
Because $G(W)\subset \mathbb H^n$, we have
\[
  G^n(z)\ge 0
\]
for all $z\in W$. Moreover,
\[
  G^n(y)=x^n=0,
\]
because $x\in\partial\mathbb H^n$. Thus $G^n$ has a local minimum at $y$.
Hence
\[
  dG^n_y=0.
\]
Therefore the image of $dG_y$ is contained in the hyperplane
\[
  \mathbb R^{n-1}\times\{0\},
\]
so
\[
  \operatorname{rank} dG_y\le n-1.
\]

On the other hand, since $F$ and $G$ are inverse diffeomorphisms, the derivative
$dG_y$ must be invertible. This is a contradiction. Hence $F(x)$ cannot be an
interior point of $\mathbb H^n$, and therefore
\[
  F(x)\in \partial\mathbb H^n.
\]
\end{proof}

\begin{lemma}
If $p\in M$ is a boundary point, then for every chart $(U,\varphi)$ containing $p$,
\[
  \varphi(p)\in \partial\mathbb H^n.
\]
\end{lemma}

\begin{proof}
Since $p$ is a boundary point, by definition there exists a chart $(V,\psi)$
containing $p$ such that
\[
  \psi(p)\in \partial\mathbb H^n.
\]
Let $(U,\varphi)$ be any chart containing $p$. On the overlap $U\cap V$, the
transition map
\[
  \varphi\circ\psi^{-1}:
  \psi(U\cap V)\to \varphi(U\cap V)
\]
is a diffeomorphism between open subsets of $\mathbb H^n$. By the invariance
of boundary for diffeomorphisms between open subsets of $\mathbb H^n$, it maps
boundary points of $\mathbb H^n$ to boundary points of $\mathbb H^n$.

Since
\[
  \psi(p)\in \partial\mathbb H^n,
\]
we obtain
\[
  \varphi(p)
  =
  (\varphi\circ\psi^{-1})(\psi(p))
  \in \partial\mathbb H^n.
\]
Therefore every chart containing $p$ sends $p$ to $\partial\mathbb H^n$.
\end{proof}
